%=============================================================================
% CROSS-REFERENCE ANALYSIS: Hypercharge Weighting w = 7/80
% How w = N_species/(g_F * Tr(Y^2)) enters the kappa_{U(1)} computation
% Generated: 2026-03-22
%=============================================================================

\documentclass[11pt]{article}
\usepackage{amsmath,amssymb,booktabs,geometry}
\geometry{margin=1in}

\begin{document}

\title{How $w = N_{\rm species}/(g_F \cdot \mathrm{Tr}(Y^2)) = 7/80$\\
Enters the $\kappa_{U(1)}$ Computation in DFD}
\author{Cross-Reference Analysis}
\date{22 March 2026}
\maketitle

%=============================================================================
\section{Summary of Findings}
%=============================================================================

After exhaustive search of all \texttt{.tex} files in v108 and the standalone
PDF ``Ab Initio Derivation of the Fine Structure Constant from Density Field
Dynamics,'' the following picture emerges.

\paragraph{Key result.}
The weighting $w = 7/80$ appears \textbf{only in the derivation-chain table}
of \texttt{section\_alpha\_locked.tex} (Table~1, line~200). It is listed as an
ingredient of the convention-locked $\alpha^{-1} = 137.03599985$ computation
but \textbf{no explicit formula of the form
$\alpha^{-1} = f(\Lambda, w, \varepsilon_{\rm adj}, \ldots)$
is written anywhere in the codebase}.

The computation proceeds through a \emph{lattice Monte Carlo / spectral-action}
pipeline, not through a single closed-form expression.  The quantity $w$
enters indirectly through the Weinberg-angle derivation and the
electroweak-mixing formula that converts $\kappa_{U(1)}$ and $\kappa_{SU(2)}$
into $\alpha$.

%=============================================================================
\section{Where $w = 7/80$ Is Defined}
%=============================================================================

\texttt{section\_alpha\_locked.tex}, Table ``Complete derivation chain for
$\alpha^{-1}$'' (lines 188--207):

\begin{center}
\renewcommand{\arraystretch}{1.2}
\begin{tabular}{llll}
\toprule
Component & Value & Source & Status \\
\midrule
$N_{\rm species}$ & $7$ & SM SU(2) components & SM content \\
$\mathrm{Tr}(Y^2)$ & $10$ & SM hypercharges & SM content \\
$g_F$ & $8$ & Spectral triple ($J \times \gamma \times \mathbb{C}$) & Derived \\
$w = N_{\rm species}/(g_F \cdot \mathrm{Tr}(Y^2))$ & $7/80$ & Hypercharge weighting & Derived \\
$\varepsilon_{\rm adj}$ & $4096/4095$ & Regular-module trace conversion & Forced \\
\bottomrule
\end{tabular}
\end{center}

%=============================================================================
\section{The Three Constituent Quantities}
%=============================================================================

\subsection{$N_{\rm species} = 7$: SM SU(2) Components}

The table labels this ``SM SU(2) components.''  Per generation, the SM contains
the following weak-isospin components (counting each SU(2) doublet component
separately and each singlet once):
\begin{itemize}
\item $Q_L = (u_L, d_L)$: 2 components (color-triplet, but SU(2)-wise = 2)
\item $L_L = (\nu_L, e_L)$: 2 components
\item $u_R$: 1 singlet
\item $d_R$: 1 singlet
\item $e_R$: 1 singlet
\end{itemize}
Total SU(2) components per generation: $2 + 2 + 1 + 1 + 1 = 7$.

This counts the number of \emph{distinct weak-isospin field species} per
generation, irrespective of color multiplicity.  (If color were included, one
would count $2\times 3 + 2 + 1\times 3 + 1\times 3 + 1 = 15$ Weyl spinors.)

\textbf{Note:} The right-handed neutrino $\nu_R$ has $Y = 0$ and does not
contribute to the hypercharge trace; its omission is consistent with
$N_{\rm species} = 7$ rather than 8.

\subsection{$\mathrm{Tr}(Y^2) = 10$}

Two different normalizations appear in the codebase:

\paragraph{Per-generation with color multiplicity (Appendix Z, PRL supplement):}
\begin{equation}
\mathrm{Tr}(Y^2) = \sum_{\text{1 gen}} d_3\, d_2\, Y^2
= 3\cdot2\cdot\tfrac{1}{36} + 3\cdot1\cdot\tfrac{4}{9}
+ 3\cdot1\cdot\tfrac{1}{9} + 1\cdot2\cdot\tfrac{1}{4}
+ 1\cdot1\cdot1
= \frac{10}{3}.
\end{equation}
This is the standard value used in GUT normalization and the Weinberg-angle
proof.

\paragraph{In the $w$-definition table:}
$\mathrm{Tr}(Y^2) = 10$ (not $10/3$).  This corresponds to summing
$\mathrm{Tr}(Y^2)$ over all three generations:
$3 \times (10/3) = 10$.

\textbf{Consistency check:}
$w = N_{\rm species}/(g_F \cdot \mathrm{Tr}(Y^2))
= 7/(8 \times 10) = 7/80$.
This is self-consistent only if $\mathrm{Tr}(Y^2)$ in the $w$-formula
is the 3-generation sum $= 10$.

\subsection{$g_F = 8$: Spectral Triple Dimension}

The table labels this ``Spectral triple ($J \times \gamma \times \mathbb{C}$) ---
Derived.'' In the Chamseddine--Connes noncommutative geometry framework, the
finite spectral triple has a KO-dimension of $6 \pmod{8}$, and the finite
Hilbert space carries a real structure $J$, a chirality $\gamma_F$, and complex
conjugation, giving $g_F = 2^3 = 8$ real degrees of freedom per complex
fermion field.

This factor rescales the spectral-action trace: when expanding $\mathrm{Tr}(f(D_A/\Lambda))$
over the finite geometry, each fermionic degree of freedom contributes
$g_F = 8$ times due to the particle/antiparticle, left/right, and
real/imaginary doublings.

%=============================================================================
\section{How $w$ Enters the $\alpha$ Computation}
%=============================================================================

\subsection{The Two-Route Structure}

The DFD codebase presents $\alpha$ through two complementary routes:

\paragraph{Route 1: Lattice Monte Carlo (Ab Initio paper, Appendix K).}
The lattice simulation measures stiffnesses $\kappa_{U(1)}$ and $\kappa_{SU(2)}$
at the derived parameter point $(\beta_{U(1)}, \beta_{SU(2)}) = (3.80, 22.80)$.
The fine-structure constant is extracted via Wilson's normalization:
\begin{equation}\label{eq:alpha-lattice}
\alpha_W = \frac{e^2}{4\pi}
= \frac{(1/\kappa_{U(1)})(4/\kappa_{SU(2)})}
       {(1/\kappa_{U(1)}) + (4/\kappa_{SU(2)})}
\cdot \frac{1}{4\pi}.
\end{equation}
In this route, $w$ does not appear explicitly---the hypercharge structure is
encoded in the lattice action's $U(1)$ sector, and the factor of 4 in
$g_2^2 = 4/\kappa_{SU(2)}$ is the standard SU(2) Wilson normalization
$\mathrm{Tr}(T^a T^b) = \tfrac{1}{2}\delta^{ab}$.

\paragraph{Route 2: Convention-locked microsector
(Section~\ref*{sec:alpha-locked}).}
This is the sub-ppm derivation.  Here the spectral action on
$\mathbb{C}P^2 \times S^3$ produces a raw gauge-kinetic coefficient from
the $a_4$ Seeley--DeWitt term.  To convert this ``raw spectral action
coefficient'' into the physical $\kappa_{U(1)}$, one must account for:

\begin{enumerate}
\item The \textbf{hypercharge trace weight} $\mathrm{Tr}(Y^2)$: how the
      spectral action distributes kinetic energy among gauge sectors.
\item The \textbf{spectral triple multiplicity} $g_F = 8$: the total
      doubling of fermionic degrees of freedom.
\item The \textbf{species count} $N_{\rm species} = 7$: how many distinct
      weak-isospin field species contribute per generation.
\item The \textbf{BOOST factor} $\varepsilon_{\rm adj} = d^2/(d^2 - 1)
      = 4096/4095$: the regular-module trace conversion.
\item The \textbf{finite-$k$ cutoff} $\Lambda^3 = k(k+3)/(k+4)$:
      the Toeplitz truncation factor.
\end{enumerate}

\subsection{The Inferred Role of $w$}

Although no single formula is written explicitly, the role of $w$ can be
reconstructed from the Chamseddine--Connes spectral action formalism.  In the
standard approach, the $a_4$ coefficient of the spectral action expansion gives:
\begin{equation}
S_{\rm gauge} \supset \frac{f_0}{2\pi^2}
\sum_r c_r \int \mathrm{Tr}(F_r^{\mu\nu} F_{r\,\mu\nu})\,d^4x,
\end{equation}
where $c_r$ are group-theoretic coefficients.  For the U(1) factor:
\begin{equation}
c_{U(1)} \propto \frac{N_{\rm species}}{g_F \cdot \mathrm{Tr}(Y^2)} = w.
\end{equation}

The physical coupling is then:
\begin{equation}
\frac{1}{g_1^2} \propto \kappa_{U(1)} = (\text{raw } a_4 \text{ coefficient})
\times w \times \varepsilon_{\rm adj} \times (\text{finite-}k \text{ factor}).
\end{equation}

In other words:
\begin{equation}
\boxed{\kappa_{U(1)} = \kappa_{\rm raw} \cdot w \cdot \varepsilon_{\rm adj}
\cdot \frac{k+3}{k+4}}
\end{equation}
where $\kappa_{\rm raw}$ is the universal gauge-kinetic prefactor from the
spectral action, and $w = 7/80$ is the \textbf{multiplicative weight that
converts the universal coefficient to the U(1) hypercharge sector}.

\subsection{Comparison with Weinberg Angle Derivation}

In the Weinberg angle proof (Appendix Z), the same hypercharge trace appears
differently.  There, the ratio of couplings is:
\begin{equation}
\frac{g'^2}{g^2}
= \frac{\kappa_{SU(2)} \cdot \tfrac{1}{2}}
       {\kappa_{U(1)} \cdot \mathrm{Tr}(Y^2)}
= \frac{2\kappa_0 \cdot \tfrac{1}{2}}{1 \cdot \kappa_0 \cdot \tfrac{10}{3}}
= \frac{3}{10},
\end{equation}
giving $\sin^2\theta_W = 3/13$.  Here $\mathrm{Tr}(Y^2) = 10/3$ is the
per-generation value, and the stiffness ratio $\kappa_{SU(2)}/\kappa_{U(1)}
= n_2/n_1 = 2$ enters.

The $w = 7/80$ formulation packages the \emph{same} physics differently:
it folds $\mathrm{Tr}(Y^2)$ (3-generation sum), $g_F$, and $N_{\rm species}$
into a single weight.

%=============================================================================
\section{Gap Analysis: What Is Missing from the Codebase}
%=============================================================================

\begin{enumerate}
\item \textbf{No explicit master formula.}  The codebase never writes
a single equation of the form
$\alpha^{-1} = F(k_{\max}, w, \varepsilon_{\rm adj}, \Lambda, \ldots)$.
The computation is distributed across:
\begin{itemize}
\item The Chern--Simons weighted sum $\to \beta_{U(1)} = 3.80$
\item The Wilson ratio $\to \beta_{SU(2)} = 22.80$
\item The lattice extraction formula (Eq.~13 in the Ab Initio PDF)
\item The BOOST factor correction
\end{itemize}

\item \textbf{$N_{\rm species} = 7$ is not derived.}  It is stated as
``SM SU(2) components'' and labeled ``SM content'' (not ``Derived'').
A derivation from the spectral triple structure showing \emph{why} one
counts 7 species (not 15 Weyl spinors, not 5 multiplet types) is absent.

\item \textbf{The $\mathrm{Tr}(Y^2) = 10$ vs.\ $10/3$ normalization}
is not explained.  The Weinberg angle proof uses $10/3$ (per generation);
the $w$-table uses $10$ (three generations).  The text does not state
which convention is used or why.

\item \textbf{$g_F = 8$ is asserted but not derived.}  The label
``Spectral triple ($J \times \gamma \times \mathbb{C}$)'' suggests
$2 \times 2 \times 2 = 8$ from particle/antiparticle $\times$
left/right $\times$ complex conjugation, but the derivation from the
finite spectral triple axioms is not given.
\end{enumerate}

%=============================================================================
\section{Answer to the Original Question}
%=============================================================================

\begin{quote}
\textbf{Q:} Is it $\kappa_{U(1)} = (\text{raw } a_4) \times w$, or
$\kappa_{U(1)} = (\text{raw } a_4) / \mathrm{Tr}(Y^2)$, or something else?
\end{quote}

\paragraph{Answer:} It is the first option, extended:
\begin{equation}
\kappa_{U(1)} = \kappa_{\rm raw} \times w \times \varepsilon_{\rm adj}
\times \frac{k+3}{k+4},
\qquad w = \frac{N_{\rm species}}{g_F \cdot \mathrm{Tr}(Y^2)} = \frac{7}{80}.
\end{equation}

The weighting $w$ is a \textbf{multiplicative rescaling factor} applied to the
universal spectral-action gauge-kinetic coefficient to obtain the U(1)
hypercharge sector's contribution.  It combines:
\begin{itemize}
\item $N_{\rm species} = 7$ (how many weak-isospin field species run in
      the U(1) loop),
\item $g_F = 8$ (spectral triple fermionic multiplicity), and
\item $\mathrm{Tr}(Y^2) = 10$ (3-generation hypercharge trace, in the
      denominator because the trace-normalization of the U(1) kinetic
      term is $\propto 1/\mathrm{Tr}(Y^2)$).
\end{itemize}

The factor $\varepsilon_{\rm adj} = 4096/4095$ is the forced BOOST from the
regular-module microsector, and $(k+3)/(k+4) = 63/64$ is the finite-$k$
Toeplitz truncation factor.

\paragraph{However:} this reconstruction is \textbf{inferred} from context.
The codebase does not contain this formula explicitly.  The actual numerical
result $\alpha^{-1} = 137.03599985$ is obtained through lattice Monte Carlo
with the sub-ppm BOOST correction applied post-hoc, not through a single
analytic formula.

%=============================================================================
\section{Source Files Referenced}
%=============================================================================

\begin{itemize}
\item \texttt{section\_alpha\_locked.tex} --- Convention-locked $\alpha$ derivation
      (lines 188--207: derivation chain table; lines 117--120: BOOST factor)
\item \texttt{appendix\_Z\_complete\_params.tex} --- Weinberg angle proof
      (lines 39--53: $\mathrm{Tr}(Y^2) = 10/3$ derivation)
\item \texttt{appendix\_K.tex} --- Microsector physics: $\alpha$ from CS theory
      and lattice verification (lines 30--35: CS weighted sum; lines 125--128:
      stiffness ratio $\kappa_{U(1)}/\kappa_{SU(2)} = 1/2$)
\item \texttt{appendix\_F.tex} --- Gauge emergence foundations
      (Theorem F.13: Frame Stiffness from Ricci curvature)
\item \texttt{supplementary/PRL\_SM\_from\_Topology.tex} --- PRL-style paper
      (lines 76--93: canonical hypercharge normalization)
\item Ab Initio PDF, Eq.~(13): lattice $\alpha$ extraction formula
\end{itemize}

\end{document}
